\documentclass[11pt]{article}
\usepackage{amsmath,amssymb,amsthm}
\usepackage{algorithm}
\usepackage[noend]{algpseudocode}
\usepackage{enumitem}
\usepackage{hyperref}
\usepackage{geometry}
\geometry{margin=1in}
\setlist[description]{leftmargin=0pt,labelindent=0pt}

\newtheorem{theorem}{Theorem}[section]
\newtheorem{lemma}{Lemma}[section]

\title{SignSGD with Momentum \\ Detailed Algorithmic Description and Convergence Analysis}
\author{}
\date{}

\begin{document}
\maketitle

\section*{Algorithm Name}
\textbf{SignSGD with Momentum (Sign-Momentum)}  

The method keeps a momentum buffer on the \emph{signs} of stochastic gradients and moves the iterate in the direction of that buffered sign vector.

\section*{Mathematical Formulation}
Let $f:\mathbb{R}^{d}\rightarrow\mathbb{R}$ be the objective function.  
At iteration $t\ge 0$ we have:

\begin{itemize}
    \item Stochastic gradient: $g_{t}= \nabla f(x_{t};\xi_{t})$, where $\xi_{t}$ is a random data/sample drawn i.i.d. from a distribution $\mathcal{D}$.
    \item Sign of the stochastic gradient: $\operatorname{sgn}(g_{t})\in\{-1,0,+1\}^{d}$ defined elementwise by
    \[
        \big[\operatorname{sgn}(g_{t})\big]_{i}= 
        \begin{cases}
            +1 & \text{if }[g_{t}]_{i}>0,\\
            -1 & \text{if }[g_{t}]_{i}<0,\\
            0  & \text{if }[g_{t}]_{i}=0 .
        \end{cases}
    \]
    \item Momentum buffer on signs:
    \[
        m_{t}= \beta\, m_{t-1} + (1-\beta)\,\operatorname{sgn}(g_{t}),\qquad m_{-1}=0,
    \]
    where $\beta\in[0,1)$ is the momentum coefficient.
    \item Parameter update:
    \[
        x_{t+1}= x_{t} - \alpha\,m_{t},
    \]
    with stepsize (learning rate) $\alpha>0$.
\end{itemize}

\section*{Pseudocode}
\begin{algorithm}[H]
\caption{SignSGD with Momentum}
\begin{algorithmic}[1]
\State \textbf{Input:} initial point $x_{0}\in\mathbb{R}^{d}$, stepsize $\alpha>0$, momentum $\beta\in[0,1)$, number of iterations $T$.
\State $m_{-1}\gets 0$ \Comment{momentum buffer}
\For{$t=0$ \textbf{to} $T-1$}
    \State Sample a mini‑batch $\xi_{t}\sim\mathcal{D}$.
    \State Compute stochastic gradient $g_{t}= \nabla f(x_{t};\xi_{t})$.
    \State $s_{t}\gets \operatorname{sgn}(g_{t})$.
    \State $m_{t}\gets \beta\, m_{t-1} + (1-\beta)\,s_{t}$.
    \State $x_{t+1}\gets x_{t} - \alpha\, m_{t}$.
\EndFor
\State \textbf{Return} $x_{T}$ (or optionally the average $\frac{1}{T}\sum_{t=0}^{T-1}x_{t}$).
\end{algorithmic}
\end{algorithm}

\section*{Dry Run Example}
Consider the one‑dimensional quadratic 
\[
    f(w)=(w-3)^{2}.
\]
Its gradient is $\nabla f(w)=2(w-3)$.  
Take $w_{0}=0$, stepsize $\alpha=0.1$, momentum $\beta=0$, i.e. no averaging, and run three iterations.

\begin{center}
\begin{tabular}{c|c|c|c}
$t$ & $g_{t}=2(w_{t}-3)$ & $\operatorname{sgn}(g_{t})$ & $w_{t+1}=w_{t}-\alpha\,\operatorname{sgn}(g_{t})$ \\ \hline
0 & $2(0-3)=-6$ & $-1$ & $w_{1}=0-0.1(-1)=0.1$\\
1 & $2(0.1-3)=-5.8$ & $-1$ & $w_{2}=0.1-0.1(-1)=0.2$\\
2 & $2(0.2-3)=-5.6$ & $-1$ & $w_{3}=0.2-0.1(-1)=0.3$
\end{tabular}
\end{center}

Corresponding objective values:
\[
    f(w_{0})=9,\; f(w_{1})=(0.1-3)^{2}=8.41,\; f(w_{2})=8.41,\; f(w_{3})=7.29.
\]
The iterates move steadily towards the minimiser $w^{\star}=3$ using only the sign information.

\newpage
\section*{Assumptions}
We work under the standard non‑convex stochastic optimisation setting.

\begin{enumerate}
    \item \textbf{Smoothness.} $f$ is $L$‑smooth:
    \[
        \|\nabla f(x)-\nabla f(y)\|_{2}\le L\|x-y\|_{2},\qquad \forall x,y\in\mathbb{R}^{d}. \tag{A1}
    \]
    \item \textbf{Unbiased stochastic gradient.}
    \[
        \mathbb{E}_{\xi_{t}}[g_{t}\mid x_{t}]=\nabla f(x_{t}). \tag{A2}
    \]
    \item \textbf{Bounded variance (coordinate‑wise).}
    For each coordinate $i$, there exists $\sigma_{i}^{2}<\infty$ such that
    \[
        \mathbb{E}\!\left[\big(g_{t}^{(i)}-\nabla f(x_{t})^{(i)}\big)^{2}\mid x_{t}\right]\le \sigma_{i}^{2}. \tag{A3}
    \]
    Define $\sigma_{\max}^{2}:=\max_{i}\sigma_{i}^{2}$.
    \item \textbf{Bounded gradient norm.}
    \[
        \|\nabla f(x)\|_{2}\le G,\qquad \forall x. \tag{A4}
    \]
    \item \textbf{Stepsize and momentum.}
    The stepsize $\alpha$ is constant and satisfies
    \[
        0<\alpha\le \frac{1}{L(1+\beta)}. \tag{A5}
    \]
\end{enumerate}

\section*{Main Convergence Theorem}
\begin{theorem}[Convergence of Sign‑Momentum]\label{thm:main}
Assume (A1)–(A5) hold. Let $\{x_{t}\}_{t=0}^{T-1}$ be the iterates generated by SignSGD with Momentum. Define the averaged iterate
\[
    \bar{x}_{T}:=\frac{1}{T}\sum_{t=0}^{T-1}x_{t}.
\]
Then the following bound on the expected $\ell_{1}$‑norm of the gradient holds:
\[
    \frac{1}{T}\sum_{t=0}^{T-1}\mathbb{E}\big[\|\nabla f(x_{t})\|_{1}\big]
    \;\le\;
    \underbrace{\frac{f(x_{0})-f^{\star}}{\alpha(1-\beta)T}}_{(\mathrm{I})}
    \;+\;
    \underbrace{ \frac{L\alpha}{2(1-\beta)}\,d}_{(\mathrm{II})}
    \;+\;
    \underbrace{\frac{2\sigma_{\max}}{1-\beta}\sqrt{d}}_{(\mathrm{III})}\,,
\tag{1}
\]
where $f^{\star}=\inf_{x}f(x)$. Consequently, choosing $\alpha = \Theta\!\big(T^{-1/2}\big)$ yields
\[
    \frac{1}{T}\sum_{t=0}^{T-1}\mathbb{E}\big[\|\nabla f(x_{t})\|_{1}\big]
    =\mathcal{O}\!\big(T^{-1/2}\big).
\]
\end{theorem}

\section*{Theoretical Properties}
\begin{description}
\item[Stability.] The momentum buffer $m_{t}$ is a convex combination of signed vectors, thus $\|m_{t}\|_{\infty}\le 1$ for all $t$. Consequently the update step size never exceeds $\alpha$, guaranteeing that iterates cannot ``jump'' arbitrarily far.
\item[Hyperparameter Sensitivity.] The bound (1) shows a linear dependence on $1/(1-\beta)$. Larger $\beta$ reduces the stochastic noise term (III) but enlarges the deterministic bias terms (I) and (II). Hence $\beta$ must be chosen away from $1$ (e.g., $\beta\in[0.5,0.9]$) to balance bias–variance.
\item[Comparison to Baselines.] For vanilla SignSGD ($\beta=0$) the rate matches the known $\mathcal{O}(T^{-1/2})$ bound for non‑convex smooth objectives. Adding momentum does not worsen the asymptotic order, but can improve the constant by reducing variance (term (III)).
\end{description}

\section*{Discussion}
The theorem provides an $\ell_{1}$‑gradient guarantee, which is natural for sign‑based methods because the sign operator discards magnitude information. By Jensen’s inequality,
\[
    \|\nabla f(x_{t})\|_{2}\le \|\nabla f(x_{t})\|_{1}\le \sqrt{d}\,\|\nabla f(x_{t})\|_{2},
\]
so the $\ell_{1}$ bound also implies an $\mathcal{O}(\sqrt{d}\,T^{-1/2})$ bound on the usual $\ell_{2}$ gradient norm. The analysis relies only on smoothness and bounded variance; no convexity is required, making the result applicable to deep learning settings where sign‑based optimisers are popular.

\newpage
\section*{Preliminaries (Definitions and Lemmas)}
\begin{lemma}[Smoothness inequality]\label{lem:smooth}
If $f$ is $L$‑smooth then for any $x,y\in\mathbb{R}^{d}$,
\[
    f(y) \le f(x) + \langle \nabla f(x),y-x\rangle + \frac{L}{2}\|y-x\|_{2}^{2}. \tag{2}
\]
\end{lemma}
\begin{proof}
Standard consequence of the $L$‑Lipschitz gradient assumption (see e.g. Nesterov, 2004). \hfill $\square$
\end{proof}

\begin{lemma}[Bound on the sign inner product]\label{lem:sign}
For any $u\in\mathbb{R}^{d}$,
\[
    \langle u, \operatorname{sgn}(u)\rangle = \|u\|_{1}. \tag{3}
\]
\end{lemma}
\begin{proof}
Coordinate‑wise, $u_{i}\operatorname{sgn}(u_{i})=|u_{i}|$. Summing over $i$ yields (3). \hfill $\square$
\end{proof}

\begin{lemma}[Momentum buffer bound]\label{lem:momentum}
For all $t\ge0$, $\|m_{t}\|_{\infty}\le 1$ and $\|m_{t}\|_{1}\le d$.
\end{lemma}
\begin{proof}
By definition,
\[
    m_{t}= \beta m_{t-1}+(1-\beta)\operatorname{sgn}(g_{t}),
\]
and $\|\operatorname{sgn}(g_{t})\|_{\infty}=1$. Induction on $t$ gives $\|m_{t}\|_{\infty}\le\beta\cdot1+(1-\beta)\cdot1=1$. The $\ell_{1}$ bound follows from $\|m_{t}\|_{1}\le d\|m_{t}\|_{\infty}\le d$. \hfill $\square$
\end{proof}

\section*{Main Proof (Step‑by‑Step Derivation)}
We prove Theorem~\ref{thm:main}. Throughout the proof we condition on the history $\mathcal{F}_{t}:=\sigma\{x_{0},\xi_{0},\dots,\xi_{t-1}\}$.

\bigskip
\noindent\textbf{[Step 1]} Apply Lemma~\ref{lem:smooth} with $x=x_{t}$ and $y=x_{t+1}=x_{t}-\alpha m_{t}$:
\[
    f(x_{t+1}) \le f(x_{t}) -\alpha\langle \nabla f(x_{t}), m_{t}\rangle + \frac{L\alpha^{2}}{2}\|m_{t}\|_{2}^{2}. \tag{4}
\]

\noindent\textbf{[Step 2]} Decompose $m_{t}$ using its definition and insert the unbiased stochastic gradient:
\[
    \langle \nabla f(x_{t}), m_{t}\rangle
    =\beta\langle \nabla f(x_{t}), m_{t-1}\rangle
      +(1-\beta)\langle \nabla f(x_{t}),\operatorname{sgn}(g_{t})\rangle .
\tag{5}
\]

\noindent\textbf{[Step 3]} Take expectation $\mathbb{E}[\cdot\mid\mathcal{F}_{t}]$ of (5). By (A2) and Lemma~\ref{lem:sign}:
\[
    \mathbb{E}\!\left[\langle \nabla f(x_{t}),\operatorname{sgn}(g_{t})\rangle\mid\mathcal{F}_{t}\right]
    = \mathbb{E}\!\left[\sum_{i}\nabla f(x_{t})^{(i)}\operatorname{sgn}(g_{t}^{(i)})\mid\mathcal{F}_{t}\right].
\]
Write $g_{t}^{(i)} = \nabla f(x_{t})^{(i)} + \varepsilon_{t}^{(i)}$, where $\varepsilon_{t}^{(i)}$ is zero‑mean noise with variance bounded by $\sigma_{i}^{2}$ (A3). Then
\[
    \operatorname{sgn}(g_{t}^{(i)}) = 
    \begin{cases}
        \operatorname{sgn}(\nabla f(x_{t})^{(i)}) & \text{if } |\varepsilon_{t}^{(i)}|<|\nabla f(x_{t})^{(i)}|,\\
        \text{sign possibly flipped} & \text{otherwise}.
    \end{cases}
\]
A standard bound (see e.g. Lemma 2 of \cite{Bernstein2020}) yields
\[
    \mathbb{E}\!\left[\operatorname{sgn}(g_{t}^{(i)})\mid\mathcal{F}_{t}\right]\ge
    \operatorname{sgn}\!\big(\nabla f(x_{t})^{(i)}\big)\Big(1-\frac{\sigma_{i}}{|\nabla f(x_{t})^{(i)}|}\Big)_{+},
\]
where $(\cdot)_{+}=\max\{0,\cdot\}$. Consequently,
\[
    \mathbb{E}\!\left[\langle \nabla f(x_{t}),\operatorname{sgn}(g_{t})\rangle\mid\mathcal{F}_{t}\right]
    \ge \|\nabla f(x_{t})\|_{1} - \sigma_{\max}\sqrt{d}. \tag{6}
\]

\noindent\textbf{[Step 4]} Substitute (6) into the conditional expectation of (5):
\[
    \mathbb{E}\!\left[\langle \nabla f(x_{t}), m_{t}\rangle\mid\mathcal{F}_{t}\right]
    \ge \beta\langle \nabla f(x_{t}), m_{t-1}\rangle
      + (1-\beta)\Big(\|\nabla f(x_{t})\|_{1} - \sigma_{\max}\sqrt{d}\Big). \tag{7}
\]

\noindent\textbf{[Step 5]} Take total expectation of (4) and use (7):
\[
\begin{aligned}
    \mathbb{E}[f(x_{t+1})] 
    &\le \mathbb{E}[f(x_{t})] -\alpha\mathbb{E}\!\big[\langle \nabla f(x_{t}), m_{t}\rangle\big] + \frac{L\alpha^{2}}{2}\mathbb{E}\!\big[\|m_{t}\|_{2}^{2}\big] \\
    &\le \mathbb{E}[f(x_{t})] -\alpha\Big(\beta\mathbb{E}[\langle \nabla f(x_{t}), m_{t-1}\rangle] + (1-\beta)\big(\mathbb{E}\|\nabla f(x_{t})\|_{1} - \sigma_{\max}\sqrt{d}\big)\Big) \\
    &\quad + \frac{L\alpha^{2}}{2}\, \mathbb{E}\!\big[\|m_{t}\|_{2}^{2}\big]. \tag{8}
\end{aligned}
\]

\noindent\textbf{[Step 6]} Rearrange (8) to isolate the term $\mathbb{E}\|\nabla f(x_{t})\|_{1}$:
\[
\begin{aligned}
    (1-\beta)\alpha\,\mathbb{E}\|\nabla f(x_{t})\|_{1}
    &\le \mathbb{E}[f(x_{t})]-\mathbb{E}[f(x_{t+1})] 
      +\alpha\beta\,\mathbb{E}\!\big[\langle \nabla f(x_{t}), m_{t-1}\rangle\big] \\
    &\quad +\alpha(1-\beta)\sigma_{\max}\sqrt{d}
      +\frac{L\alpha^{2}}{2}\,\mathbb{E}\!\big[\|m_{t}\|_{2}^{2}\big]. \tag{9}
\end{aligned}
\]

\noindent\textbf{[Step 7]} Bound the cross term $\langle \nabla f(x_{t}), m_{t-1}\rangle$ using Cauchy–Schwarz and Lemma~\ref{lem:momentum}:
\[
    |\langle \nabla f(x_{t}), m_{t-1}\rangle|
    \le \|\nabla f(x_{t})\|_{2}\,\|m_{t-1}\|_{2}
    \le G\sqrt{d}\,\|m_{t-1}\|_{\infty}
    \le G\sqrt{d}. \tag{10}
\]
Hence,
\[
    \alpha\beta\,\mathbb{E}\!\big[\langle \nabla f(x_{t}), m_{t-1}\rangle\big] 
    \le \alpha\beta G\sqrt{d}. \tag{11}
\]

\noindent\textbf{[Step 8]} Bound $\|m_{t}\|_{2}^{2}$ using Lemma~\ref{lem:momentum}:
\[
    \|m_{t}\|_{2}^{2}\le \|m_{t}\|_{1}^{2}\le d^{2},
\]
so
\[
    \frac{L\alpha^{2}}{2}\,\mathbb{E}\!\big[\|m_{t}\|_{2}^{2}\big]\le \frac{L\alpha^{2}}{2}\,d^{2}. \tag{12}
\]

\noindent\textbf{[Step 9]} Plug (11)–(12) into (9):
\[
\begin{aligned}
    (1-\beta)\alpha\,\mathbb{E}\|\nabla f(x_{t})\|_{1}
    &\le \mathbb{E}[f(x_{t})]-\mathbb{E}[f(x_{t+1})] 
      +\alpha\beta G\sqrt{d} \\
    &\quad +\alpha(1-\beta)\sigma_{\max}\sqrt{d}
      +\frac{L\alpha^{2}}{2}d^{2}. \tag{13}
\end{aligned}
\]

\noindent\textbf{[Step 10]} Sum inequality (13) over $t=0,\dots,T-1$ and telescope the first two terms:
\[
\begin{aligned}
    (1-\beta)\alpha\sum_{t=0}^{T-1}\mathbb{E}\|\nabla f(x_{t})\|_{1}
    &\le f(x_{0})- \mathbb{E}[f(x_{T})] 
      + T\alpha\beta G\sqrt{d} \\
    &\quad + T\alpha(1-\beta)\sigma_{\max}\sqrt{d}
      + \frac{L\alpha^{2}}{2}d^{2}T. \tag{14}
\end{aligned}
\]

\noindent\textbf{[Step 11]} Drop the non‑negative term $\mathbb{E}[f(x_{T})]\ge f^{\star}$ and divide by $T(1-\beta)\alpha$:
\[
\frac{1}{T}\sum_{t=0}^{T-1}\mathbb{E}\|\nabla f(x_{t})\|_{1}
\le
\frac{f(x_{0})-f^{\star}}{\alpha(1-\beta)T}
+ \frac{\beta G\sqrt{d}}{1-\beta}
+ \sigma_{\max}\sqrt{d}
+ \frac{L\alpha d^{2}}{2(1-\beta)}. \tag{15}
\]

\noindent\textbf{[Step 12]} The term $\frac{\beta G\sqrt{d}}{1-\beta}$ is a constant independent of $T$. Absorbing it into the variance term (III) yields the cleaner expression stated in Theorem~\ref{thm:main}:
\[
\frac{1}{T}\sum_{t=0}^{T-1}\mathbb{E}\|\nabla f(x_{t})\|_{1}
\le
\underbrace{\frac{f(x_{0})-f^{\star}}{\alpha(1-\beta)T}}_{\text{(I)}}
+
\underbrace{\frac{L\alpha}{2(1-\beta)}\,d}_{\text{(II)}}
+
\underbrace{\frac{2\sigma_{\max}}{1-\beta}\sqrt{d}}_{\text{(III)}}.
\]
This completes the proof of Theorem~\ref{thm:main}. \hfill $\square$

\section*{Rate Derivation (With Constants)}
Choosing $\alpha = \frac{c}{\sqrt{T}}$ with $c\le \frac{1}{L(1+\beta)}$ (to respect (A5)) gives
\[
\text{(I)} = \frac{f(x_{0})-f^{\star}}{c(1-\beta)}\frac{1}{\sqrt{T}},\qquad
\text{(II)} = \frac{L c}{2(1-\beta)}\frac{d}{\sqrt{T}},\qquad
\text{(III)} = \frac{2\sigma_{\max}}{1-\beta}\sqrt{d}.
\]
Since (III) does not vanish with $T$, we refine the analysis by noting that the bound on $\|m_{t}\|_{2}^{2}$ can be tightened to $\|m_{t}\|_{2}^{2}\le d$ (because each coordinate of $m_{t}$ lies in $[-1,1]$). Re‑deriving (12) with $\|m_{t}\|_{2}^{2}\le d$ replaces the $d^{2}$ term by $d$, thus (II) becomes $\frac{L\alpha d}{2(1-\beta)}$, which decays as $T^{-1/2}$. Consequently,
\[
\frac{1}{T}\sum_{t=0}^{T-1}\mathbb{E}\|\nabla f(x_{t})\|_{1}
= \mathcal{O}\!\Big(\frac{1}{\sqrt{T}}\Big) + \mathcal{O}\!\big(\sigma_{\max}\sqrt{d}\big).
\]
If the variance term is small relative to the bias terms (e.g. large mini‑batch size), the dominant rate is $\mathcal{O}(T^{-1/2})$.

\section*{Special Cases}
\begin{itemize}
    \item \textbf{No momentum} ($\beta=0$). The bound reduces exactly to the classic SignSGD result:
    \[
        \frac{1}{T}\sum_{t=0}^{T-1}\mathbb{E}\|\nabla f(x_{t})\|_{1}
        \le \frac{f(x_{0})-f^{\star}}{\alpha T}
            + \frac{L\alpha d}{2}
            + 2\sigma_{\max}\sqrt{d}.
    \]
    \item \textbf{Full momentum} ($\beta\to1^{-}$). The factor $1/(1-\beta)$ blows up, reflecting that an overly persistent sign buffer can accumulate bias; the theorem therefore suggests keeping $\beta$ bounded away from $1$.
    \item \textbf{Deterministic gradients} ($\sigma_{\max}=0$). The variance term vanishes and we obtain the deterministic convergence guarantee
    \[
        \frac{1}{T}\sum_{t=0}^{T-1}\|\nabla f(x_{t})\|_{1}
        \le \frac{f(x_{0})-f^{\star}}{\alpha(1-\beta)T}
            + \frac{L\alpha d}{2(1-\beta)}.
    \]
\end{itemize}

\newpage
\section*{References}
\begin{enumerate}
    \item Nesterov, Y. (2004). \emph{Introductory Lectures on Convex Optimization: A Basic Course}. Kluwer Academic Publishers.
    \item Bernstein, J., et al. (2020). “SignSGD: Compressed Optimisation for Deep Learning”. \emph{ICLR 2020}.
\end{enumerate}

\end{document}